\section{Dismissal, blessing, and the last Gospel (nn.~1132--1139)}

The last eight numbers of the \work{Ordo Missae} are the shortest stretch of
the book and the most legislated. Three of the four things they contain ---
the dismissal formula, the blessing, and the last Gospel --- are conditional,
and the conditions are spelled out both in the \work{Ordo} and in the general
rubrics. A reader who wants to see how the 1960 code works can do worse than
read these eight numbers against nn.~507--510.

\subsection*{The dismissal (nn.~1132--1136)}

\begin{appointed}{1132--1133 $\cdot$ \latincue{Ite, missa est}}
\textit{Focused source locator.} The appointed unit is identified by this incipit and marginal number; only phrases required by the analysis below are quoted.
\end{appointed}

\begin{englishwitness}{xliv}
``Ite missa est (vel) Benedicamus Domino. Go, you are dismissed, (or) let us
bless the Lord. \textbf{R.} Deo gratias. \textbf{R.} Thanks be to God.'' And,
in Masses for the dead: ``\textbf{P.} Requiescant in pace. \textbf{P.} May they
rest in peace. \textbf{R.} Amen.''
\end{englishwitness}

\begin{rubricnote}{\work{Rubricae generales Missalis romani} nn.~507--508}
\latincue{Ite, missa est} is said at the end of Mass and answered \latincue{Deo
gratias}. But: at the evening Mass of Maundy Thursday followed by the solemn
reposition, and at other Masses followed by a procession, \latincue{Benedicamus
Domino} is said instead; within the octave of Easter, in Masses of the season,
a double \latincue{Alleluia} is added to \latincue{Ite, missa est} and to the
answering \latincue{Deo gratias}; and at Requiems \latincue{Requiescant in
pace} is said, answered \latincue{Amen}. The blessing is given after
\latincue{Placeat}, and is omitted only where \latincue{Benedicamus Domino} or
\latincue{Requiescant in pace} has been said.
\end{rubricnote}

The formula is famously hard to construe, and the 1861 translator's ``Go, you
are dismissed'' is one of the two defensible readings. \latincue{Missa} here is
a late-Latin noun for a dismissal, and \latincue{Ite, missa est} says
approximately: go, the dismissal is made. What has to be resisted is the
temptation to read \latincue{missa} in the sentence as the noun it later became
--- the Mass --- since the name of the rite is derived from this sentence and
not the other way round. The 1962 book offers no gloss and this study offers no
resolution beyond noting which way the derivation runs.

Three things about the choice among the three formulas deserve notice, because
each is a small piece of ecclesiology. \latincue{Benedicamus Domino} replaces
the dismissal exactly when nobody is being dismissed: when a procession
follows, the assembly is not sent away but taken somewhere. \latincue{Requiescant
in pace} replaces it at a Requiem, so that the last words addressed to the
congregation at a funeral Mass are not addressed to the congregation at all. And
the answer to the ordinary form is \latincue{Deo gratias}, which is the same
answer given to the Epistle: the Mass ends as its lessons end.

The missal then prints eight chant settings of \latincue{Ite, missa est} under
n.~1134, keyed to the rank and season of the day in the same way as the
settings of the \latincue{Gloria}: one with a double \latincue{Alleluia} for
Easter week, one for the rest of paschal time, and six more for first-, second-
and third-class feasts, feasts of Our Lady, the Sundays of the year, and the
plainest class of day. Two settings of \latincue{Benedicamus Domino} follow at
n.~1135 and one of \latincue{Requiescant in pace} at n.~1136. The last and
simplest of the eight carries a rubric permitting it \latincue{ad libitum} at
any Mass in which \latincue{Ite, missa est} is sung without
\latincue{Alleluia}.

\subsection*{\latincue{Placeat} and the blessing (nn.~1137--1138)}

\begin{appointed}{1137--1138 $\cdot$ \latincue{Placeat tibi} and \latincue{Benedicat vos}}
\textit{Focused source locator.} The appointed unit is identified by this incipit and marginal number; only phrases required by the analysis below are quoted.
\end{appointed}

\begin{englishwitness}{xliv}
``Let the performance of my homage be pleasing to thee, O holy Trinity; and
grant that the sacrifice which I, tho' unworthy, have offered up in the sight
of thy Majesty, may be acceptable to thee, and thro' thy mercy be a propitiation
for me, and all those for whom it has been offered. Thro' \dots'' And the
blessing: ``May Almighty God, the Father, Son, and Holy Ghost, bless you.
Amen.''
\end{englishwitness}

\latincue{Placeat} is the priest's last private prayer and it says something
the public prayers of the Mass never say in the first person singular: that the
sacrifice was offered by an unworthy man, and that its acceptability and its
propitiatory effect are still being asked for after it has been offered and
consumed. \latincue{Obsequium servitutis meae} --- the homage of my service ---
is a formula of subjection, and the 1861 ``the performance of my homage''
catches its flatness. The prayer is said bowed, with joined hands on the altar,
in the same posture as \latincue{Supplices te rogamus}.

The blessing is one blessing, and the missal makes a point of saying so:
\latincue{semel tantum benedicens, etiam in Missis solemnibus} --- blessing only
once, even at solemn Masses. The single exception is printed immediately
after: \latincue{In Missa pontificali ter benedicitur, ut in Pontificali
habetur}, at a pontifical Mass the blessing is threefold, as the
\work{Pontificale} has it. The formula is divided across the turn: the priest
says \latincue{Benedicat vos omnipotens Deus} facing the altar, then turns to
the people for the three names. The grammar is one sentence and the ceremony
splits it in two, so that the naming of the Trinity is the part said to the
congregation.

\subsection*{The last Gospel (n.~1139)}

\begin{appointed}{1139 $\cdot$ \latincue{Initium sancti Evangelii secundum Ioannem} (Ioann.~1, 1--14)}
\textit{Focused source locator.} The appointed unit is identified by this incipit and marginal number; only phrases required by the analysis below are quoted.
\end{appointed}

\begin{englishwitness}{xlv--xlvi}
``IN the beginning was the Word, and the Word was with God, and the Word was
God \dots\ And THE WORD WAS MADE FLESH, and dwelt among us; and we saw his
glory, as it were the glory of the only begotten of the Father, full of grace
and truth. \textbf{R.} Thanks be to God.'' The 1861 book sets \latincue{ET
VERBUM CARO FACTUM EST} and ``AND THE WORD WAS MADE FLESH'' in capitals, as it
had set \latincue{ET HOMO FACTUS EST} in the Creed, and closes the Ordinary
with a printer's rule immediately after \latincue{Deo gratias}.
\end{englishwitness}

\begin{rubricnote}{When it is not John, and when there is none: nn.~509--510}
n.~509: for the last Gospel, in any Mass, the beginning of the Gospel according
to John is regularly taken. However, on the second Sunday of the Passion or
Palm Sunday, in all Masses that do not follow the blessing and procession of
palms, a proper last Gospel is said. n.~510: the last Gospel is omitted
altogether in six cases --- in Masses in which \latincue{Benedicamus Domino}
has been said under n.~507 a; on Christmas Day at the third Mass; on Palm
Sunday, in the Mass that follows the blessing and procession of palms; at the
Mass of the Easter Vigil; at Requiems when the absolution at the catafalque
follows; and in certain Masses following consecrations, from the rubrics of the
\work{Pontificale Romanum}.
\end{rubricnote}

Two of those six are worth pausing on because they show the rubric reasoning.
At the third Mass of Christmas the last Gospel is omitted because John 1,
1--14 has just been read as the Gospel of the day: the rite will not read the
same passage twice in one Mass. On Palm Sunday, in the Mass that follows the
procession, the last Gospel is omitted for the opposite reason --- the rite has
already been long --- while in every other Mass of the same day a proper last
Gospel, and not John, is appointed. The general principle is that the last
Gospel is the first thing to go.

That is a fact about its status, and the fact has a history.

\begin{witnesscard}{Historical note: how John 1 became part of the Mass}
Fortescue's account is that throughout the Middle Ages the last Gospel belonged
to the priest's thanksgiving and ``was in no wise an element of the Mass''. The
Sarum Missal has it said on the way back to the sacristy; Burchard's ceremonial
of 1502 allows it to be said at the altar; and ``Pius V in his reformed missal
(1570) for the first time admits it as part of the Mass''. Fortescue adds the
observations that make the point stick: it is not sung by the deacon, it has no
solemnities at High Mass, and a bishop says it on his way from the
altar.\footnote{Fortescue, \work{The Mass}, p.~394, read at the page image.}
This is an identified early twentieth-century reconstruction, reported as such.
What can be checked in the 1962 book itself agrees with it: the last Gospel is
the one reading in the Mass with no chant, no ministers' part beyond two
responses, and a list of six occasions on which it simply does not happen.
\end{witnesscard}

The genuflection at \latincue{Et Verbum caro factum est} recalls the Creed's
kneeling clause. At \latincue{Et incarnatus est} the Incarnation is professed;
here John narrates it. The consecration stands between those two textual
moments, but it makes Christ sacramentally present and must not be described as
a repetition of the Incarnation. Whether any compiler intended a three-part
symbolic sequence cannot be shown from the book, and this study claims none.

The Mass ends with \latincue{Deo gratias}, the answer given to the Epistle and
to the dismissal. The 1962 \work{Ordo Missae} closes there, at n.~1139, and the
next number in the volume is the introit of Easter Sunday.
